% !TeX root = euclid-without-words.tex
% !TeX Program=pdfLaTeX
%
% Inscribed angle one half of central angle
%
\begin{minipage}[t]{.48\textwidth}
\hspace{3em}
\begin{tikzpicture}[scale=.8]
% Circle goes from (a) to (b)
\coordinate (a) at (0,0);
\coordinate (b) at (6,0);
\coordinate (center) at (3,0);
% Line containing lower points is (c) -- (d)
\coordinate (c) at (0,-2);
\coordinate (d) at (6,-2);
% Line containing upper points is (e) -- (f);
\coordinate (e) at (0,2.3);
\coordinate (f) at (4.5,3);
% Name the upper and lower lines
\path [name path=chord1] (c) -- (d);
\path [name path=chord2] (e) -- (f);
% Draw a circle whose center is half-way between (a) and (b) through (a)
\node [draw,thick,circle through=(a),name path=circ] at (center) {};
% Get intersections of the upper and lower lines with the circle
\path [name intersections={of=circ and chord1,by={i1,i2}}];
\path [name intersections={of=circ and chord2,by={i3,i4}}];
% Draw the lower chord
%\draw[thick] (i1) -- (i2);
% Draw the two subtended angles and the center angle
\draw[thick,red] (i1) -- (i3)
  node[below,xshift=-3pt,yshift=-12pt] {$\bm{\alpha}$} -- (i2);
\draw[thick,black!10!teal] (i1) -- (i4)
  node[below,xshift=6pt,yshift=-12pt] {$\bm{\alpha}$} -- (i2);
\draw[thick,blue] (i1) -- (center)
  node[below,yshift=-8pt] {$\bm{2\alpha}$} -- (i2);
\end{tikzpicture}
\end{minipage}
